\section{The Indian setting}\label{sec:setting}

This section describes the rules that shape each limit to arbitrage in Indian stock futures. \Cref{fig:timeline} places
the dated changes on one time line.

\begin{figure}[t]
\caption{Almost every limit to arbitrage changed on a known date during the sample}\label{fig:timeline}
\exsub{Futures STT, enforcement of the ban, the position-limit rule and arbitrage capital, 2015 to 2026}
\centering\includegraphics[width=\linewidth]{fig_timeline.pdf}
\exsource{Source: SEBI circulars and press releases; NSE circulars (listed in the references).}
\end{figure}

\subsection{Stock futures}
Single-stock futures trade on the NSE for about two hundred large and liquid stocks, the members of the F\&O segment. A
stock enters the segment when it meets eligibility criteria set by the Securities and Exchange Board of India (SEBI),
based on market capitalisation, traded value and delivery, and leaves it when it fails them. Contracts expire monthly,
three months are listed at a time, and most trading is in the nearest month. Since October 2019 stock derivatives are
settled by delivery of the shares. One detail of the exchange's prices matters for measurement: when a future does not
trade in the last half hour of a day, the exchange sets its settlement price at a theoretical cost-of-carry value. A
basis computed from settlement prices would therefore look small on thin days for purely mechanical reasons, so I use
prices from actual trades throughout.

\subsection{Short selling}
SEBI's framework allows all investors to sell short, but only in securities of the F\&O segment, and it forbids naked
short selling: every short sale must be delivered at settlement, so an overnight short needs borrowed shares. Shares can
be borrowed through the exchange's securities lending and borrowing (SLB) segment, but it is thin: on a typical day only
a few dozen stocks are lent. In practice the overnight short position in an Indian stock is a short future. This is why
the reverse arbitrage trade (short the stock, buy the future) is much costlier than cash-and-carry, and why F\&O
membership defines which stocks can be shorted at all.

\subsection{The position limit and the ban}
The MWPL caps the total open interest that all traders together may hold in a stock's futures and options. Until
September 2025 it was 20\% of the shares held by non-promoters (the free float), and open interest was counted gross,
with every option at its full notional value. From 1 October 2025 the limit is the lower of 15\% of the free float and
65 times the average daily delivery volume, with a floor of 10\% of the free float, recomputed every quarter, and open
interest is counted in delta-adjusted (futures-equivalent) terms (\cref{tab:mwpl}). The new formula lowered every
stock's limit, to a median of 73\% of its old value.

\begin{table}[t]
\caption{The limit was cut and open interest redefined in October 2025}\label{tab:mwpl}
\exsub{Market-wide position limit (MWPL) and the open interest tested against it}
\small\begin{tabularx}{\textwidth}{@{}K{3.2cm}LL@{}}
\toprule
\thd{Period} & \thd{MWPL} & \thd{Open interest tested} \\
\midrule
To 30 Sep 2025 & 20\% of shares held by non-promoters (free float) & Gross open interest, futures and options \\ \rowrule
23 Mar to 26 Nov 2020 & Halved for volatile or heavily used stocks & Gross open interest \\ \rowrule
From 1 Oct 2025 & Lower of 15\% of free float and 65 times average daily delivery (in shares), floor 10\% of free float; recomputed quarterly & Futures-equivalent (delta-adjusted) open interest \\
\bottomrule
\end{tabularx}
\exsource{Source: SEBI circulars of 22 September 2006 and 29 May 2025; SEBI press release 18/2020; NSE Clearing circular 119/2025.}
\end{table}

The ratio of open interest to the limit is the stock's \emph{utilisation}. When utilisation exceeds 95\% at the end of a
day, the stock is in the ban from the next day; the ban ends after utilisation falls to 80\%. During the ban every
participant may only reduce positions, in all expiries and in both futures and options, and rolling a position into the
next month counts as a new position. \Cref{fig:sail} shows the rule at work for SAIL, the stock banned most often.

\begin{figure}[t]
\caption{A stock enters the ban when open interest crosses 95\% of its limit and leaves at 80\%}\label{fig:sail}
\exsub{Utilisation of the market-wide position limit for SAIL, January to June 2024; shaded days are ban days}
\centering\includegraphics[width=\linewidth]{fig_util_example.pdf}
\exsource{Source: NSE combined open interest files and securities-in-ban lists; author's calculations.}
\end{figure}

How violations were detected changed twice, and this turns out to matter. Until 22 March 2020 an increase was judged on
end-of-day positions, so a position opened and closed within the day was not a breach. From 23 March 2020, following
SEBI's press release 18/2020, the clearing corporations judged violations on peak intraday positions. From 1 October
2025 each entity's end-of-day delta-adjusted position must fall relative to its position when the ban began, which
allows new contracts that lower the delta, such as a put bought against a long future.

\subsection{Trading costs and arbitrage capital}
The STT on the sale of futures was 0.01\% of the traded value from June 2013, 0.0125\% from 1 April 2023, 0.02\% from 1
October 2024 and 0.05\% from 1 April 2026. A cash-and-carry position pays much more on the stock: 0.1\% STT on the
purchase and 0.1\% on delivery at expiry, plus stamp duty, exchange charges and brokerage. The main holders of
cash-and-carry positions are arbitrage mutual funds, which are taxed as equity funds. When debt funds lost their tax
advantage on 1 April 2023, short-term money moved into arbitrage funds, and their assets grew from Rs 71,106 crore in
April 2023 to Rs 2.98 lakh crore in September 2026.
